\section{Aktualisierung der Dateneinfügungsskripte (Data Seeding und Ingestion)}

Durch die Einführung des Blind-Indexing-Verfahrens mussten sämtliche bestehenden Migrations-, Import- und Ingestion-Skripte grundlegend überarbeitet werden. Da die Datenbank nun für jedes verschlüsselte Attribut zwingend den korrespondierenden Blind Index benötigt, können Rohdaten nicht mehr nativ oder über klassische SQL-Inserts direkt in die Tabellen eingepflegt werden.

Die Skripte wurden dahingehend modifiziert, dass jeder Datensatz vor dem Schreibvorgang eine kryptografische Pipeline auf Applikationsebene durchläuft. Folgende Prozessschritte wurden in den automatisierten Einfügungsskripten implementiert:

\begin{enumerate}
    \item \textbf{Klartext-Validierung:} Die einzufügenden Rohdaten (z.\,B. aus CSV- oder JSON-Quellen) werden auf Validität und korrekte ISO-Formatierung (insbesondere Datums- und Telefonnummernformate) geprüft.
    \item \textbf{Generierung des Blind Index:} Für die betroffenen Felder (\texttt{firstName}, \texttt{lastName}, \texttt{plz}, \texttt{city}, \texttt{street}, \texttt{birthday}, \texttt{phone}, \texttt{email}) wird deterministisch ein HMAC-SHA256-Hash unter Verwendung des applikationsspezifischen geheimen Schlüssels (\textit{Secret Pepper}) berechnet und Base64-kodiert im jeweiligen \texttt{Hash}-Feld hinterlegt.
    \item \textbf{Probabilistische Verschlüsselung:} Der eigentliche Nutzinhalt wird mittels eines sicheren, authentifizierten Verschlüsselungsverfahrens (AES-GCM-256) mit einem pro Datensatz einzigartigen, zufälligen Initialisierungsvektor verschlüsselt.
    \item \textbf{Persistierung:} Der finale \texttt{INSERT}-Befehl schreibt sowohl den unvorhersehbaren Chiffretext als auch den indizierbaren Blind Index synchron in die Datenbank.
\end{enumerate}

%\begin{figure}[h]
%    \centering
%    % Hier kann bei Bedarf ein Flowchart-Pfad hinterlegt werden
%    \caption{Pipeline der Dateneinfügungsskripte für PII-Ressourcen.}
%    \label{fig:data_ingestion_pipeline}
%\end{figure}

Die Skripte stellen somit sicher, dass die referenzielle Integrität und die Suchbarkeit der verschlüsselten Personenobjekte (\texttt{Patients}, \texttt{Doctors}, \texttt{Nurses}) ab dem exakten Zeitpunkt der Dateninjektion konsistent gewahrt bleiben.